% Source printed pp. 504-506 / PDF pp. 512-514.
\section{簇的模空间}\label{src:10.1}
我们详细介绍带标点稳定曲线的模空间 $\overline M_{g,n}$ 的若干自然推广.

\subsection{加权带标点稳定曲线}\label{src:10.1.1-subsection}
Hassett 的\newterm{加权带标点稳定曲线}{weighted pointed stable curve}模空间 $\overline M_{g,A}$ 自然推广了 $\overline M_{g,n}$, 允许某些标记点重合 \cite{Has03}. 它也为 \S\bookref{1.7} 中概述的射影模空间六步构造法提供了典型的例证.

固定亏格 $g$ 和权向量 $A=(a_1,\ldots,a_n)\in\mathbb Q^n$, 满足对所有 $i$ 有 $0<a_i\leq1$, 且 $2g-2+\sum a_i>0$. 一条\newterm{$A$-加权稳定曲线}{$A$-weighted stable curve}是亏格为 $g$ 的连通结点射影曲线 $C$, 带有标记点 $p_1,\ldots,p_n\in C$, 满足每个 $p_i\in C$ 是光滑点, 对每个 $q\in C$ 有 $\sum_{p_i=q}a_i\leq1$, 且 $\mathbb Q$-线丛 $\omega_C(\sum a_ip_i)$ 丰沛. $A$-加权稳定曲线族是带有 $n$ 个截面的平坦固有有限表现态射, 其几何纤维满足上述条件. 一条 $(1,\ldots,1)$-加权稳定曲线恰好就是带标点稳定曲线 (定义~\bookref{6.3.2}).

参数化 $A$-加权稳定曲线的模叠 $\mathcal M_{g,A}$ 是 $\Spec\mathbb Z$ 上的光滑固有 Deligne--Mumford 叠, 具有射影粗模空间 $\overline M_{g,A}$. 光滑性由带标点稳定曲线的形变理论 (命题~\bookref{6.3.26}) 推广到加权稳定曲线而得; 固有性可利用赋值判别法, 从 $\overline{\mathcal M}_{g,n}$ 的固有性 (定理~\bookref{6.5.23}) 推得; 射影性则可用 Koll\'ar 丰沛性判别法 (\bookref{6.9.2}) 证明.

若权向量 $A=(a_1,\ldots,a_n)$ 与 $B=(b_1,\ldots,b_n)$ 对所有 $i$ 满足 $a_i\geq b_i$, 则有约化映射 $\overline{\mathcal M}_{g,A}\to\overline{\mathcal M}_{g,B}$. 若子集 $A'\subseteq A$ 满足 $2g-2+\sum_{a'\in A'}a'>0$, 则有遗忘映射 $\overline{\mathcal M}_{g,A}\to\overline{\mathcal M}_{g,A'}$. 另一个引人注目的特点是 $\mathbb R^n$ 中的壁与室分解: 每个室的内部, 模空间 $\overline M_{g,A}$ 保持不变. 选取特定的权向量 $A$, 可得到 $M_{g,n}$ 的其他紧化, 包括 Kapranov, Keel 和 Losev--Manin 先前构造的紧化. 例如, 若 $A=(2/3,2/3,2/3,\epsilon,\ldots,\epsilon)$, 其中 $0<\epsilon\ll1$, 则 $\overline M_{0,A}=(\mathbb P^1)^{n-3}$.

\subsection{稳定映射}\label{src:10.1.2-subsection}
Kontsevich 引入了到光滑射影簇 $X$ 的\newterm{稳定映射}{stable map}的模叠 $\overline{\mathcal M}_{g,n}(X,\beta)$, 其中类 $\beta\in\mathrm{CH}_1(X)$ \cite{Kon95}. 稳定映射由预稳定曲线 $(C,p_i)$ (即带有互异光滑点 $p_i\in C$ 的结点曲线 $C$) 以及满足 $f_*[C]=\beta$ 的映射 $f:C\to X$ 组成, 并要求: 对每个被 $f$ 收缩为一点的 $E\cong\mathbb P^1\subseteq C$ (相应地, 算术亏格为 $1$ 的 $E\subseteq C$), $E$ 至少含有 $3$ 个 (相应地, $1$ 个) 特殊点. 当 $X=\mathbb P^n$ 时, 最后这个条件等价于 $\omega_C(p_1+\cdots+p_n)\otimes f^*\mathcal O_{\mathbb P^n}(3)$ 丰沛. 稳定映射族定义为预稳定族 $(\mathcal C\to S,\sigma_i)$ 连同映射 $f:\mathcal C\to S$, 其几何纤维满足上述条件.

稳定映射叠 $\overline{\mathcal M}_{g,n}(X,\beta)$ 是对角态射有限的固有代数叠, 具有射影粗模空间 $\overline M_{g,n}(X,\beta)$. 在特征 $p$ 下, $\overline{\mathcal M}_{g,n}(X,\beta)$ 可能不是 Deligne--Mumford 叠, 因为可能存在具有 $\mu_p$-稳定子的稳定映射, 如 $\mathbb P^1\to\mathbb P^1$, $[x:y]\mapsto[x^p:y^p]$. 可按 \S\bookref{1.7} 的六步策略构造模空间 $\overline M_{g,n}(X,\beta)$, 并用 Koll\'ar 丰沛性判别法验证射影性. 详细介绍见 \cite{FP97}.

\newterm{Gromov--Witten 不变量}{Gromov--Witten invariant}由稳定映射的模空间定义. 对每个 $j$, 有取值映射 $\mathrm{ev}_j:\overline{\mathcal M}_{g,n}(X,\beta)\to X$, 定义为 $(C,p_i,f)\mapsto f(p_j)$. 给定上同调类 $\gamma_1,\ldots,\gamma_n\in\mathrm{CH}^*(X)$, 相应的 Gromov--Witten 不变量为
\[
\langle\gamma_1,\ldots,\gamma_n\rangle^X_{g,n,\beta}
:=\int_{[\overline{\mathcal M}_{g,n}(X,\beta)]^{\mathrm{virt}}}
\mathrm{ev}_1^*\gamma_1\cup\cdots\cup\mathrm{ev}_n^*\gamma_n,
\]
其中 $[\overline{\mathcal M}_{g,n}(X,\beta)]^{\mathrm{virt}}\in\mathrm{CH}_{(1-g)(\dim X-3)-\beta(\omega_X)+n}(\overline{\mathcal M}_{g,n}(x))$ 是一个``\newterm{虚基本类}{virtual fundamental class}'' \cite{LT98,BF97}. 若 $g=0$, 则 $\overline{\mathcal M}_{0,n}(X)$ 光滑且具有预期维数, 虚基本类与通常的基本类一致. $X$ 中曲线的模空间还有其他紧化, 能产生相关的计数不变量. 例如, 参数化 Calabi--Yau 三维簇中曲线的理想层的固有模空间给出 Donaldson--Thomas 不变量 \cite{DT98}; 猜想认为它们与 Gromov--Witten 不变量相关 \cite{MNOP06}.

稳定映射的模空间可推广为允许靶是固有 Deligne--Mumford 叠 $\mathcal X$ \cite{AV02}, 或者具有有限对角态射的固有驯服代数叠 $\mathcal X$ \cite{AOV08}. 若只考虑从普通曲线 $C$ 出发的稳定映射 $C\to\mathcal X$, 模空间并不紧. 因此改为允许 $C$ 是\newterm{扭曲稳定曲线}{twisted stable curve}, 即某种叠式曲线 (见例~\bookref{4.9.16}), 其结点和标记点可具有 $\mu_n$-稳定子. 稳定映射的这种推广可构造出有趣的簇的模空间; 例如取 $\mathcal X=\overline{\mathcal M}_{1,1}$, 就得到椭圆纤维化的模空间. 目前正有人尝试将理论推广到更一般的代数叠中的稳定映射: 取 $\mathcal X=B\mathbb G_m$ (相应地, $B\mathrm{GL}_n$), 将得到 $\overline{\mathcal M}_g$ 上万有 Picard 簇的紧化, 类似于 \cite{Cap94} (相应地, $\overline{\mathcal M}_g$ 上半稳定向量丛的万有模空间, 类似于 \cite{Gie84,Pan96,NS99}).
